\chapter{Trace Anomalies and Background Variations}
\label{ch:curved-trace-anomalies}

\ConventionsEuclidean

The stress tensor of Chapter~\ref{ch:point-splitting-stress-tensor-renormalization}
is defined by variation with respect to a background metric.  This chapter
studies the trace obtained by Weyl variation.  The central object is the
renormalized generating functional as a functional of background fields.

\section{Generating Functional and Weyl Variation}

Let \(W[g]\) be the renormalized connected Euclidean generating functional of
a QFT on a smooth \(D\)-manifold with metric \(g_{\mu\nu}\).  The stress
tensor expectation value is defined by
\[
  \delta W[g]
  =
  \frac12\int_M d^Dx\,\sqrt g\,
  \langle T^{\mu\nu}(x)\rangle_g\,\delta g_{\mu\nu}(x)
\]
for metric variations in the chosen renormalization scheme.  A local Weyl
variation is
\[
  \delta_\sigma g_{\mu\nu}=2\sigma g_{\mu\nu}.
\]
Therefore
\[
  \delta_\sigma W[g]
  =
  \int_M d^Dx\,\sqrt g\,
  \sigma(x)\langle T^\mu{}_\mu(x)\rangle_g .
\]
If the classical theory is Weyl invariant, a nonzero local functional on the
right side is a trace anomaly.

\section{Locality and Wess--Zumino Consistency}

Assume the regulator and renormalization prescription produce a local Weyl
variation
\[
  \delta_\sigma W[g]=\int_M d^Dx\,\sqrt g\,\sigma\,\mathcal A[g].
\]
Since Weyl transformations commute,
\[
  [\delta_{\sigma_1},\delta_{\sigma_2}]W=0.
\]
This gives the Wess--Zumino consistency condition
\[
  \delta_{\sigma_1}\!\int\sqrt g\,\sigma_2\mathcal A
  =
  \delta_{\sigma_2}\!\int\sqrt g\,\sigma_1\mathcal A .
\]
An anomaly is therefore a cohomology class of local Weyl variations modulo
Weyl variations of local counterterms.  If
\[
  W[g]\mapsto W[g]+\int_M\sqrt g\,L_{\rm ct}(g)
\]
is a finite local counterterm change, then \(\mathcal A\) changes by the Weyl
variation of \(L_{\rm ct}\).

\section{Two Dimensions}

In two dimensions, for a unitary conformal theory on a curved background, the
metric trace anomaly has the form
\[
  \langle T^\mu{}_\mu\rangle_g
  =
  -\frac{c}{24\pi}R
\]
with the sign tied to the convention for \(W[g]\) and the Euclidean action.
The coefficient \(c\) is the same parameter that appears in the flat-space
stress-tensor two-point function after the normalization of \(T_{\mu\nu}\) is
fixed.  The integral of \(R\) is topological on a closed surface, but the
local relation is stronger: it controls the response to a space-dependent
Weyl factor.

\section{Four Dimensions}

In four dimensions, the parity-even trace anomaly of a conformal theory has
the form
\[
  \langle T^\mu{}_\mu\rangle_g
  =
  \frac{1}{16\pi^2}
  \left(c\,W_{\mu\nu\rho\sigma}W^{\mu\nu\rho\sigma}
  -a\,E_4
  +b\,\Box R\right)
\]
plus background gauge-field and parity-odd terms when such backgrounds are
present.  Here
\[
  E_4=
  R_{\mu\nu\rho\sigma}R^{\mu\nu\rho\sigma}
  -4R_{\mu\nu}R^{\mu\nu}+R^2
\]
is the Euler density and \(W_{\mu\nu\rho\sigma}\) is the Weyl tensor.  The
\(\Box R\) coefficient is scheme dependent because it can be shifted by a
local \(R^2\) counterterm.  The coefficients \(a\) and \(c\) are invariant
under such local metric counterterm changes in ordinary four-dimensional CFT
settings.

\section{What Must Be Proved in a QFT}

The displayed classifications use locality of the renormalized background
variation.  In a concrete QFT one must still prove or assume that:
\begin{enumerate}[label=(\roman*)]
\item \(W[g]\) exists in the stated class of backgrounds;
\item metric variations define renormalized stress-tensor insertions;
\item Weyl variation is local after the allowed counterterms are chosen;
\item contact terms in multiple stress-tensor insertions are part of the same
  renormalization scheme;
\item anomalies in other background fields are compatible with Weyl
  consistency.
\end{enumerate}
These conditions are automatic neither in an abstract local-net formulation
nor in a symbolic path integral.  They are extra structure connecting the QFT
to background geometry.

\begin{openproblem}[Trace anomaly from nonperturbative QFT data]
\label{op:trace-anomaly-nonperturbative-data}
Formulate trace anomalies directly from nonperturbative local QFT data on
curved backgrounds.  The construction should identify the algebraic or
functorial object replacing \(W[g]\), define stress-tensor contact terms, and
recover the local Weyl cohomology class without relying on a perturbative
regulator.
\end{openproblem}
